\chapter{Classical Spaces and Group Identification}
\label{chap:classical_group}

This chapter records the standalone \texttt{ClassicalGroup} formalization
ported from \texttt{inducedorbittoy}.  It is independent of the PBP/MYD
combinatorics in the rest of this blueprint: the Lean module path is
\texttt{CombUnipotent.ClassicalGroup}, while the Lean namespace remains
\texttt{ClassicalGroup}.

The full natural-language source is kept in
\texttt{lean/docs/blueprints/classical\_group/blueprint.md}; the formalization
plan and problem statement are stored in the same directory.

\section{Core Structures}

\begin{definition}[Classical family]
  \label{def:classical_star}
  \lean{ClassicalGroup.ClassicalStar}
  \leanok
  The six families are \(B,D,C,\widetilde C,C^*,D^*\).  Each family carries
  signs \(\epsilon\), \(\dot\epsilon\), and \(\epsilon\dot\epsilon\) governing
  the bilinear form, the square of \(L\), and the square of \(J\).
\end{definition}

\begin{definition}[Formed complex space]
  \label{def:formed_space}
  \lean{ClassicalGroup.FormedSpace}
  \leanok
  \uses{def:classical_star}
  A formed space is a finite-dimensional complex vector space equipped with a
  non-degenerate \(\epsilon\)-symmetric complex bilinear form.
\end{definition}

\begin{definition}[\(J\)- and \(L\)-structures]
  \label{def:jl_structures}
  \lean{ClassicalGroup.JStructure, ClassicalGroup.LStructure}
  \leanok
  \uses{def:formed_space}
  A \(J\)-structure is a conjugate-linear form-preserving automorphism with the
  prescribed square.  An \(L\)-structure is a complex-linear form-preserving
  automorphism with the prescribed square.
\end{definition}

\section{Existence and Uniqueness}

\begin{definition}[Classical-space existence package]
  \label{def:has_classical_space}
  \lean{ClassicalGroup.HasClassicalSpace}
  \leanok
  \uses{def:classical_star, def:formed_space, def:jl_structures}
  The predicate \texttt{HasClassicalSpace star p q} asserts the existence of a
  formed complex space with compatible \(J\) and \(L\) satisfying the dimension,
  positivity, square, commutation, and signature conditions of the
  classical-space theorem.
\end{definition}

\begin{theorem}[Existence of classical spaces]
  \label{thm:exists_classical_space}
  \lean{ClassicalGroup.exists_classical_space}
  \leanok
  \uses{def:has_classical_space}
  Every legal classical signature \((\star,p,q)\) has a classical space.
\end{theorem}

\begin{definition}[Uniqueness package]
  \label{def:unique_classical_spaces}
  \lean{ClassicalGroup.UniqueClassicalSpacesFor}
  \leanok
  \uses{def:has_classical_space}
  The uniqueness predicate states that any two classical spaces for the same
  signature are related by a complex-linear form isomorphism intertwining both
  \(J\) and \(L\).
\end{definition}

\begin{theorem}[Classical-space existence and uniqueness]
  \label{thm:classical_space_existence_and_uniqueness}
  \lean{ClassicalGroup.classical_space_existence_and_uniqueness}
  \leanok
  \uses{thm:exists_classical_space, def:unique_classical_spaces}
  The core linear-algebra theorem gives existence and uniqueness for every
  legal classical signature, before adding the group-identification layer.
\end{theorem}

\section{Group Identification Layer}

\begin{definition}[Centralizer of \(J\)]
  \label{def:complex_form_isometry_centralizer_j}
  \lean{ClassicalGroup.ComplexFormIsometryCentralizerJ}
  \leanok
  \uses{def:formed_space, def:jl_structures}
  The centralizer of \(J\) is represented as the type of complex form isometries
  commuting with \(J\).  A subgroup version gives it an honest group object.
\end{definition}

\begin{definition}[Expected real group identification]
  \label{def:has_expected_group_identification}
  \lean{ClassicalGroup.HasExpectedGroupIdentification}
  \leanok
  \uses{def:complex_form_isometry_centralizer_j, def:classical_star}
  This abstract predicate records property (10) of the source theorem: the
  centralizer of \(J\) should identify with the expected real classical group.
  The current formalization keeps the concrete real Lie group API isolated.
\end{definition}

\begin{theorem}[Full classical-group theorem]
  \label{thm:classical_group_theorem}
  \lean{ClassicalGroup.classical_group_theorem}
  \leanok
  \uses{thm:classical_space_existence_and_uniqueness, def:has_expected_group_identification}
  For every legal classical signature, the Lean theorem packages the
  existence/uniqueness theorem together with the abstract group-identification
  layer.
\end{theorem}
